\section{The Frozen Simulation Study \frozenmark}
\label{sec:frozen}

\subsection{Design and freeze}

Eighty closed-loop runs were executed before any physical run: four
calibration levels, two planners, two obstacle geometries
(\textsf{K0}, obstacle on the approach line; \textsf{K1}, offset
\SI{0.35}{\metre}), five repetitions per cell. The scenario constants
shared by both domains -- nominal surge \SI{0.15}{\metre\per\second},
engagement range \SI{1.5}{\metre}, assumed obstacle height
\SI{0.5}{\metre} -- are read by both the simulated campaign and the
physical pipeline from a single file, so the two cannot drift apart
silently. The aggregate was hashed and committed; we verified during the
final audit that all 80 runs still reproduce the committed aggregate
exactly, and that every run carries the sentinel of its own calibration
level with no leakage between levels.

\subsection{What the levels do to closed-loop behaviour}

Table~\ref{tab:frozen} reports the median over five repetitions for each
cell. Three patterns are worth stating.

\emph{The ladder is not monotone in safety.} Minimum clearance for the
committed planner in \textsf{K0} rises from \SI{1.48}{\metre} at S0 to
\SI{2.22}{\metre} at S1, then falls to \SI{0.77}{\metre} at S2 and
\SI{0.58}{\metre} at S3. Adding a measured observation model made the
simulated vehicle \emph{safer}, because a biased range estimate that
reads short triggers engagement earlier; adding measured timing then
removed most of that margin, because the burst structure delays
qualification until the obstacle is close. A ladder that simply
degraded performance step by step would have been easier to report and
would have told us less.

\emph{Qualification, not detection, is what the timing model costs.} The
range at the first planner-valid observation collapses from about
\SI{3.4}{\metre} at S0 and S1 to about \SI{1.2}{\metre} at S2 and S3 in
every cell. The vehicle is not blind at long range; it is unable to
accumulate enough coherent evidence to be allowed to act.

\emph{The two planners respond to calibration in opposite directions.}
The committed planner's manoeuvre shortens from \SI{65}{\second} to
\SI{51}{\second} between S0 and S3; the dynamic-window planner's
lengthens from \SI{60}{\second} to \SI{92}{\second}. This was the
sharpest testable claim in the frozen set, and it does not depend on
absolute agreement with anything.

\begin{table}[t]
\centering
\caption{Frozen simulated campaign \frozenmark, median of five
repetitions, for the three quantities the text uses. \emph{clearance}:
minimum obstacle clearance (\si{\metre}); \emph{span}: manoeuvre span
(\si{\second}); \emph{qualified}: range at the first qualified
observation (\si{\metre}). Lateral excursion and commitment distance are
in the accompanying repository; the dynamic-window commitment distance
equals the starting distance at every level, so it is a left-censored
bound rather than a measurement.}
\label{tab:frozen}
\setlength{\tabcolsep}{5pt}
\begin{tabular}{llrrrrrrrrrrrr}
\toprule
& & \multicolumn{4}{c}{clearance} & \multicolumn{4}{c}{span}
& \multicolumn{4}{c}{qualified} \\
\cmidrule(lr){3-6}\cmidrule(lr){7-10}\cmidrule(lr){11-14}
geo & planner & S0 & S1 & S2 & S3 & S0 & S1 & S2 & S3
& S0 & S1 & S2 & S3 \\
\midrule
K0 & committed & 1.48 & 2.22 & 0.77 & 0.58
   & 65 & 56 & 56 & 51 & 3.43 & 3.03 & 1.15 & 1.17 \\
K0 & dwa & 2.14 & 2.47 & 0.93 & 0.86
   & 60 & 86 & 89 & 92 & 3.47 & 3.43 & 1.23 & 1.21 \\
K1 & committed & 1.19 & 2.56 & 0.87 & 0.74
   & 57 & 55 & 56 & 49 & 3.45 & 3.21 & 1.18 & 1.17 \\
K1 & dwa & 2.46 & 2.33 & 0.87 & 0.65
   & 73 & 31 & 90 & 92 & 3.50 & 3.09 & 1.21 & 1.21 \\
\bottomrule
\end{tabular}
\end{table}

\subsection{Failures inside the frozen set}

One run of eighty ends in a collision (S2, \textsf{K0}, committed;
clearance \SI{0.22}{\metre}). Four dynamic-window runs -- two at S0, two
at S3 -- have no commitment distance at all, so those cells are medians
of four; the same left-censoring that afflicts the physical runs
(Section~\ref{sec:realresults}) is already present here, for the same
reason, in a planner that begins manoeuvring immediately. Two S2 runs
never delivered a single planner-valid observation.

\subsection{Reporting the estimator work}

An earlier phase of this project compared four temporal estimator
variants under the same qualification. We keep that result only in the
compressed form it deserves here: the fixed-gain filter used throughout
this paper was selected on that basis, and the choice is held constant
in every run reported, so no result below depends on it. Readers
interested in the comparison will find it in the accompanying
repository; giving it more room would compete with the paper's one
story.
